\section{Introduction}


These are notes on Perelman's papers
``The Entropy Formula for the Ricci Flow and its Geometric
Applications'' \cite{Perelman} and ``Ricci Flow with Surgery
on Three-Manifolds' \cite{Perelman2}.
In these two remarkable preprints, which were posted on the ArXiv
in  2002 and 2003, Grisha Perelman
announced a proof of the Poincar\'e Conjecture, and more generally Thurston's
Geometrization Conjecture, using the Ricci flow approach of Hamilton.
Perelman's proofs are concise and,
at times, sketchy. The purpose of these notes is to provide the details
that are missing in
\cite{Perelman} and \cite{Perelman2}, which contain Perelman's arguments for
the Geometrization Conjecture.

Among other things,
we cover the construction of the Ricci flow with surgery
of \cite{Perelman2}.
We also discuss the long-time behavior of the
Ricci flow with surgery, which is needed for the full Geometrization
Conjecture.
The papers \cite{Colding-Minicozzi,Perelman3},
which are not covered in these notes,
each provide a shortcut in the case of the Poincar\'e Conjecture.
Namely, these papers show
that if the initial manifold is simply-connected then the
Ricci flow with surgery becomes extinct in a
finite time, thereby removing the issue of the long-time
behavior. Combining this claim with the proof of existence of
Ricci flow with surgery gives
the shortened proof in the simply-connected case.

These notes are intended for readers with a solid background in geometric analysis.
Good sources for background material on Ricci flow are
\cite{Chow-Knopf,Chow-Lu-Ni,Hamilton,Topping}.
The notes are self-contained but are designed to be read along with
\cite{Perelman,Perelman2}.
For the most part we follow the format of
\cite{Perelman,Perelman2} and use the section numbers
of \cite{Perelman,Perelman2} to label our
sections. We have done this in order to respect the
structure of \cite{Perelman,Perelman2} and to facilitate the use of
the present notes as a companion to \cite{Perelman,Perelman2}.
 In some places we have rearranged Perelman's arguments
or provided alternative arguments,
but we have refrained from an overall reorganization.

Besides providing details for Perelman's proofs, we have included some
expository material in the form of overviews and appendices.
Section \ref{overview0} contains an overview of the Ricci flow approach to
geometrization of $3$-manifolds. Sections \ref{overview1} and \ref{overview2} contain
overviews of \cite{Perelman} and \cite{Perelman2}, respectively.
The appendices discuss some background material and
techniques that are used throughout the notes.

Regarding the proofs, the papers \cite{Perelman,Perelman2}  contain some incorrect statements
and incomplete arguments, which we have attempted to point out to the reader.
(Some of the mistakes in \cite{Perelman} were corrected in \cite{Perelman2}.)
We did not find any
serious problems, meaning problems that cannot be corrected using the methods
introduced by Perelman.

We will refer to Section X.Y of \cite{Perelman} as I.X.Y, and
Section X.Y of \cite{Perelman2} as II.X.Y.
A reader may wish to start with the overviews, which
explain the logical structure of the arguments and the
interrelations between the sections.
It may also be helpful to browse through the appendices before delving into the
main body of the material.


These notes have gone through various versions, which were
posted at \cite{KleinerLottwebsite}.
An initial version with notes on \cite{Perelman}
was posted in June 2003.
A version covering \cite{Perelman,Perelman2} was posted in September
2004.
After the May 2006 version of these notes
 was  posted on the ArXiv, expositions of Perelman's
work appeared in
\cite{Cao-Zhu} and \cite{Morgan-Tian}.
 \\ \\
\noindent
{\em Acknowledgements.}
In the preparation of the September 2004 version of our notes,
we benefited from a
workshop on Perelman's surgery procedure
that was held in August 2004, at Princeton.
We thank the participants of that meeting, as well as
the Clay Mathematics Institute for supporting it.

We are grateful for comments and corrections that we have
received regarding earlier versions of these
notes.
We especially thank G\'erard Besson for comments on Theorem
\ref{noncollthm} and Lemma \ref{bdrysmoothness}, Peng Lu for
comments on Lemma \ref{II.4.5}, Bernhard Leeb for
discussions on the issue of uniqueness for the standard
solution and Richard Bamler for comments on
Proposition \ref{propII.6.3} and Proposition \ref{propII.6.4}.
Along with these people we thank
Mike Anderson, Albert Chau,
Ben Chow, Xianzhe Dai, Jonathan Dinkelbach,
Sylvain Maillot, John Morgan,
Joan Porti, Gang Tian, Peter Topping, Bing Wang,
Guofang Wei, Hartmut Weiss,
Burkhard Wilking, Jon Wolfson, Rugang Ye and Maciej Zworski.


We thank the referees for their detailed and helpful comments on
this paper.


We thank the Clay Mathematics Institute for supporting the writing
of the notes on \cite{Perelman2}.


\tableofcontents


\section{A reading guide}

Perelman's papers contain a wealth of results about Ricci flow.
We cover all of these results, whether or not they
are directly relevant to the Poincar\'e and Geometrization Conjectures.

Some readers may wish to take an abbreviated route that
focuses on the proof of the
Poincar\'e Conjecture or the Geometrization Conjecture.
Such readers can try the following itinerary.

Begin with the overviews in Sections \ref{overview0} and \ref{overview1}.
Then review Hamilton's compactness theorem and its variants, as described in
Appendix \ref{subsequence}; an exposition is in
\cite[Chapter 7]{Topping}.
Next, read I.7 (Sections \ref{I.7}-\ref{I.7.3}),
followed by I.8.3(b) (Section \ref{secI.8.3}).
After reviewing the theory of Riemannian
manifolds and Alexandrov spaces of nonnegative sectional curvature
(Appendix \ref{Alexandrov} and references therein),
proceed to I.11 (Sections \ref{secI.11.1}-\ref{secI.11.9}), followed by
II.1.2 and I.12.1 (Sections \ref{II.1.2}-\ref{I.12.1}).


At this point, the reader should be ready for
the overview of Perelman's second paper in Section \ref{overview2},  and
can proceed with II.1-II.5 (Sections
\ref{notationterminology}-\ref{surgeryflow}).  In conjunction with
one of the finite extinction time results
\cite{Colding-Minicozzi,Colding-Minicozzi2,Perelman3},
this completes the proof of the Poincar\'e Conjecture.

To proceed with the rest of the proof of the Geometrization Conjecture,
the reader can begin with the large-time estimates for nonsingular
Ricci flows, which appear in I.12.2-I.12.4
(Sections \ref{I.12.2}-\ref{I.12.4}).
The reader can then go to II.6 and II.7 (Sections \ref{II.6}-\ref{II.7.4}).

The main topics that are missed by such an abbreviated route are
the ${\mathcal F}$ and ${\mathcal W}$ functionals
(Sections \ref{I.1.1}-\ref{I.5}), Perelman's differential Harnack inequality
(Section \ref{I.9}),  pseudolocality (Sections
\ref{statement of I.10.1}-\ref{I.10.5}) and Perelman's alternative proof
of cusp incompressibility (Section \ref{II.8}).




\section{An overview of the Ricci flow approach to
$3$-manifold geometrization}
\label{overview0}

This section is an overview of the Ricci flow approach to
$3$-manifold geometrization.  We make no attempt to present the
history of the ideas that go into the argument.
We caution the reader that for the sake of readability,
in many places we have suppressed technical points and deliberately
oversimplified the story.
 The overview will introduce the argument in
three passes, with successively greater precision and detail : we start with a
very crude sketch, then expand this to a step-by-step outline of the
strategy, and then move on to more detailed commentary on specific
points.

Other overviews may be found in \cite{Cao-Chow,Morgan}.  The
primary objective of our exposition is
to prepare the reader for a more detailed study of Perelman's work.


We refer the reader to Appendix \ref{geom} for the statement
of the geometrization conjecture.



By convention, all manifolds and Riemannian metrics in this section
will be smooth.
We follow the notation of
\cite{Perelman,Perelman2} for pointwise quantities:  $R$
denotes the scalar curvature,
$\Ric$ the Ricci curvature,  and $|\Rm|$ the largest absolute value of
the sectional curvatures.
An inequality such as $\Rm\geq C$ means that all of the
sectional curvatures at a
point or in a region, depending on the context, are bounded below by $C$.
In this section, we will specialize to three dimensions.


\subsection{The definition of Ricci flow, and some basic properties}


Let $M$ be a compact  $3$-manifold and let $\{g(t)\}_{t\in [a,b]}$
be a smoothly varying family of  Riemannian metrics on $M$.
Then $g(\cdot)$ satisfies the {\em Ricci flow equation} if
\begin{equation}
\frac{\D g}{\D t}(t) \: = \:
- \: 2 \: \Ric(g(t))
\end{equation}
holds for every $t\in[a,b]$.
Hamilton showed in \cite{Hamiltonnnnn} that
for any Riemannian metric $g_0$ on $M$, there is a $T\in (0,\infty]$
with the property that there is a (unique) solution $g(\cdot)$
to the Ricci flow equation defined on the time interval $[0,T)$ with
$g(0) = g_0$, so
that if $T<\infty$ then the curvature of $g(t)$ becomes unbounded
as $t\ra T$.  We refer to this maximal solution as the Ricci flow with
initial condition $g_0$. If $T<\infty$ then we call $T$ the blow-up
time.  A basic example
is the shrinking  round $3$-sphere, with $g_0 \: = \: r_0^2 \: g_{S^3}$ and
$g(t) \: = \: (r_0^2 - 4t) \:  g_{S^3}$, in which case $T \: = \: \frac{r_0^2}{4}$.

Suppose that $M$ is simply-connected.
Based on the round $3$-sphere example, one could hope that every Ricci flow
on $M$ blows up in finite time and becomes round while shrinking to a  point,
as $t$ approaches the blow-up time $T$. If so, then
by rescaling and taking a limit as $t \rightarrow T$, one would show
that $M$ admits a metric of constant positive sectional curvature and
therefore,
by a classical theorem, is diffeomorphic to $S^3$. The analogous argument does work
in two dimensions \cite[Chapter 5]{Chow-Knopf}. Furthermore, if the initial
metric $g_0$
 has positive Ricci curvature then Hamilton
showed in \cite{Hamilton} that this is the correct picture:
the manifold shrinks to a point in finite time and becomes round
as it shrinks.

One is then led to ask what can happen if $M$ is simply-connected
but $g_0$ does not have positive Ricci curvature.  Here a new
phenomenon can occur --- the Ricci flow solution may become singular
before it has time to shrink to a point, due to a possible neckpinch.
A neckpinch is modeled by a product region $(- c, c) \times S^2$ in which
one or many $S^2$-fibers separately shrink to a point at time $T$, due to the
positive curvature of $S^2$.   The formation of  neckpinch (and other)
singularities prevents one from continuing the Ricci flow.
In order to continue the evolution some intervention is required, and this
is the role of surgery.  Roughly speaking, the
idea of surgery is to remove a neighborhood diffeomorphic to
$(- c^\prime, c^\prime) \times S^2$ containing the shrinking $2$-spheres,
and cap off the resulting boundary components by
gluing in $3$-balls.   Of course the
topology of the manifold changes during surgery -- for instance it may become
disconnected -- but it changes in a controlled way.
The postsurgery Riemannian manifold is smooth,
so one may  restart the Ricci flow using
it as an initial condition.
When continuing the flow one may encounter further neckpinches,
which give rise to further surgeries, etc.
One hopes that
eventually all of the connected components
shrink to points while becoming round, i.e. that the Ricci flow solution
has a {\em finite extinction time}.


\subsection{A rough outline of the Ricci flow proof of the Poincar\'e
Conjecture}
We now give a step-by-step glimpse of the proof,
stating the needed steps as claims.

One starts with a compact orientable $3$-manifold $M$ with an
arbitrary metric $g_0$. For the moment we do not assume that
$M$ is simply-connected. Let $g(\cdot)$ be the Ricci flow
with initial condition $g_0$, defined on $[0,T)$.
Suppose that $T <\infty$.   Let $\Omega \subset M$ be
the set of points $x \in M$ for which
$\lim_{t \rightarrow T^-} R(x,t) $ exists
and is finite.
Then $M - \Omega$ is the part of $M$ that is going singular.
(For example, in the case of a single standard neckpinch, $M - \Omega$
is a $2$-sphere.) The first claim says what $M$ looks like
near this singularity set.

\begin{proposition}[existence for Riemannian metrics]
  \label{prop:kl-overview-ricci-flow-approach-manifold-geometrization-existence-riemannian-metrics}
  \dcref{Claim 1}
  \group{Kleiner--Lott Overview}
  \source{kleiner-lott}
 \cite{Perelman2} \label{introclaim1}
The set $\Omega$ is  open and as $t\ra T$, the evolving metric
$g(\cdot)$ converges smoothly
on compact subsets of $\Omega$ to
a Riemannian metric $\overline{g}$.
There is a geometrically defined neighborhood
$U$ of $M - \Omega$ such that each connected component of $U$
is either \\ \\
\noindent
A. Compact and diffeomorphic to $S^1 \times S^2$,
$S^1 \times_{\Z_2} S^2$ or $S^3/\Gamma$, where $\Gamma$ is a finite
subgroup of $\SO(4)$ that acts freely and isometrically on the round $S^3$.
(In writing $S^1 \times_{\Z_2} S^2$, the generator of $\Z_2$ acts on
$S^1$ by complex conjugation and on $S^2$ by the antipodal map.
Then $S^1 \times_{\Z_2} S^2$ is diffeomorphic to $\R P^3 \# \R P^3$.)

or \\ \\
\noindent
B. Noncompact and diffeomorphic to $\R\times S^2$,
$\R^3$
or the twisted line
bundle $\R \times_{\Z_2} S^2$ over $\R P^2$.  \\

In Case B, the connected component meets $\Omega$ in
geometrically controlled collar regions diffeomorphic to $\R\times S^2$.
\end{proposition}

Thus Claim \ref{introclaim1} provides a topological description of
a neighborhood $U$ of the region
$M-\Omega$ where the Ricci flow is going singular, along with some
geometric control on $U$.


\begin{proposition}[existence for Ricci flow with surgery]
  \label{prop:kl-overview-ricci-flow-approach-manifold-geometrization-existence-ricci-flow-surgery}
  \dcref{Claim 2}
  \group{Kleiner--Lott Overview}
  \source{kleiner-lott}
 \cite{Perelman2} \label{introclaim2}
There is a well-defined way
to perform surgery on $M$, which yields a smooth  post-surgery
manifold $M'$ with a Riemannian metric $g'$.
\end{proposition}

Claim \ref{introclaim2} means that there is a well-defined
procedure for  specifying the part of $M$ that will be removed, and
for gluing caps on the resulting manifold with boundary.
The discarded part corresponds to the neighborhood $U$ in Claim \ref{introclaim1}.
The procedure is required to satisfy a number
of additional conditions which we do not mention here.

Undoing the surgery, i.e. going from the postsurgery manifold to the
presurgery manifold, amounts to restoring some discarded
components (as in Case A of Claim \ref{introclaim1})
and performing connected sums of some
of the components of the postsurgery manifold, along with some
possible connected sums with a finite number of new $S^1 \times S^2$ and
$\R P^3$ factors.  The $S^1 \times S^2$ comes from the case when a
surgery does not disconnect the connected component where it is
performed.  The $\R P^3$ factors arise from the twisted line bundle
components in Case B of Claim \ref{introclaim1}.

After performing a surgery one lets the new manifold evolve under
the Ricci flow until one encounters the next blowup time (if
there is one).  One then performs further surgery, lets the new manifold
evolve, and repeats the process.

\begin{proposition}[Ricci flow with surgery proposition]
  \label{prop:kl-overview-ricci-flow-approach-manifold-geometrization-ricci-flow-surgery-proposition}
  \dcref{Claim 3}
  \group{Kleiner--Lott Overview}
  \source{kleiner-lott}
 \cite{Perelman2} \label{introclaim3}
One can arrange the surgery procedure so that the surgery times
do not accumulate.
\end{proposition}

If the surgery times were to accumulate, then one would have trouble continuing
the flow further, effectively killing the whole program.
Claim \ref{introclaim3} implies that by alternating Ricci flow and surgery,
one  obtains an evolutionary process that is defined for all time
(though the manifold may become the empty set from some time
onward).  We call this   {\em Ricci flow with surgery}.


\begin{proposition}[Ricci flow with surgery proposition]
  \label{prop:kl-overview-ricci-flow-approach-manifold-geometrization-ricci-flow-surgery-proposition-4}
  \dcref{Claim 4}
  \group{Kleiner--Lott Overview}
  \source{kleiner-lott}
 \cite{Colding-Minicozzi,Colding-Minicozzi2,Perelman3}
\label{introclaim4}
If the original manifold $M$ is simply-connected then any Ricci flow with
surgery on $M$ becomes extinct in finite time.
\end{proposition}

Having a finite extinction time means that from some time onwards, the
manifold is the empty set.
More generally, the same proof shows that
if the prime decomposition of the original manifold $M$
has no aspherical factors, then every Ricci flow with surgery on $M$
becomes extinct in finite time.
(Recall that a connected manifold $X$ is aspherical if
$\pi_k(X) = 0$ for all $k > 1$ or, equivalently, if its
universal cover is contractible.)

The Poincar\'e Conjecture follows immediately from the above claims.
From Claim \ref{introclaim4},
after some finite time the manifold is the empty set.
From Claims \ref{introclaim1}, \ref{introclaim2} and \ref{introclaim3},
the original manifold $M$ is diffeomorphic to a connected sum of
factors that are each $S^1 \times S^2$ or a standard quotient
$S^3/\Gamma$ of $S^3$. As we are assuming that $M$ is simply-connected,
van Kampen's theorem implies that $M$ is diffeomorphic to a connected
sum of $S^3$'s, and hence is diffeomorphic to $S^3$.





\subsection{Outline of the proof of the Geometrization Conjecture}
We now drop the assumption that $M$ is
simply-connected. The main difference is that
Claim \ref{introclaim4} no longer applies, so the Ricci flow with surgery may
go on forever in a nontrivial way.  (We remark that Claim \ref{introclaim4}
is needed only for a shortened proof of the Poincar\'e Conjecture; the proof
in the general case is logically independent of Claim \ref{introclaim4} and
also implies the Poincar\'e Conjecture.)
The possibility that there are infinitely many surgery times is
not excluded, although it is not
known whether this can actually happen.

A simple example of a Ricci flow that does not become extinct is when
$M = H^3/\Gamma$, where $\Gamma$ is a freely-acting cocompact discrete subgroup
of the orientation-preserving isometries of hyperbolic space $H^3$. If
$g_{hyp}$ denotes the metric on $M$ of constant sectional curvature
$-1$ and $g_0 \: = \: r_0^2 \: g_{hyp}$ then
$g(t) \: = \: (r_0^2 \: + \: 4t) \: g_{hyp}$. Putting
$\widehat{g}(t) \: = \: \frac{1}{t} \: g(t)$, one finds that
$\lim_{t \rightarrow \infty} \widehat{g}(t)$ is the metric on $M$
of constant sectional curvature $- \: \frac14$, independent of
$r_0$.

Returning to the general case,
let $M_t$ denote the time-$t$ manifold in a Ricci flow with surgery.
(If $t$ is a surgery time
then we consider the postsurgery manifold.)
If for some $t$
a component of
$M_t$ admits a metric with
nonnegative scalar curvature then one can show  that
the component becomes extinct or admits a flat metric; either possiblity
is good enough when we are trying to prove the Geometrization
Conjecture for the initial manifold $M$.
So we will assume
that for every  $t$, each component of $M_t$ has a point with strictly
negative scalar curvature.


Motivated by the hyperbolic
example, we consider the metric $\widehat{g}(t) \: = \: \frac{1}{t} \: g(t)$
on $M_t$. Given $x \in M_t$, define the {\em intrinsic scale} $\rho(x,t)$ to be
the radius $\rho$ such that $\inf_{B(x,\rho)} \Rm \: = \: - \: \rho^{-2}$,
where $\Rm$ denotes the sectional curvature of $\widehat{g}(t)$; this
is well-defined because the scalar curvature is negative somewhere in
the connected component of $M_t$ containing $x$.
Given $w > 0$, define the
{\em $w$-thick part} of $M_t$ by
\begin{equation}
M^+(w,t) \: = \:
\{ x \in M_t \: : \: \vol(B(x, \rho(x,t))) \: > \: w \: \rho(x,t)^3 \}.
\end{equation}
It is not excluded that $M^+(w,t) \: = \: M_t$ or $M^+(w,t) \: = \:
\emptyset$.
The next claim says that for any $w > 0$, as time goes on,
$M^+(w,t)$ approaches the $w$-thick part of a manifold of
constant sectional curvature $- \frac14$.

\begin{proposition}[existence for sectional curvature]
  \label{prop:kl-overview-ricci-flow-approach-manifold-geometrization-existence-sectional-curvature}
  \dcref{Claim 5}
  \group{Kleiner--Lott Overview}
  \source{kleiner-lott}
 \cite{Perelman2} \label{introclaim5}
There is a finite collection
$\{ (H_i, x_i) \}_{i=1}^k$ of complete pointed finite-volume $3$-manifolds
with constant sectional curvature $- \frac14$ and, for large $t$,
a decreasing function $\alpha(t)$ tending to zero and a family
of maps
\begin{equation}
f_t \: : \bigsqcup_{i=1}^k H_i\:\supset\: \bigsqcup_{i=1}^k
\:\:B \left( x_i, \frac{1}{\alpha(t)} \right)
 \rightarrow M_t\,,
\end{equation}
such that\\
1. $f_t$ is $\alpha(t)$-close to being an isometry. \\
2. The image of $f_t$ contains $M^+(\alpha(t), t)$.\\
3. The image under $f_t$ of a cuspidal torus of
$\{H_i\}_{i=1}^k$ is incompressible in $M_t$.
\end{proposition}

The proof of Claim \ref{introclaim5} uses earlier work by
Hamilton \cite{Hamiltonn}.

\begin{proposition}[graph manifolds proposition]
  \label{prop:kl-overview-ricci-flow-approach-manifold-geometrization-graph-manifolds-proposition}
  \dcref{Claim 6}
  \group{Kleiner--Lott Overview}
  \source{kleiner-lott}
 \cite{Perelman2,Shioya-Yamaguchi} \label{introclaim6}
Let $Y_t$ be the truncation of $\bigcup_{i=1}^k H_i$ obtained by
removing horoballs at distance approximately $\frac{1}{2 \alpha(t)}$
from the basepoints $x_i$. Then for large $t$, $M_t - f_t(Y_t)$ is
a graph manifold.
\end{proposition}

Claim \ref{introclaim6} reduces to a statement in Riemannian geometry
about $3$-manifolds that are locally volume-collapsed with a lower
bound on sectional curvature.

Claims \ref{introclaim5} and \ref{introclaim6}, along with Claims
\ref{introclaim1}-\ref{introclaim3}, imply the geometrization conjecture,
cf. Appendix \ref{geom}.

In the remainder of this section, we will discuss some of the claims in
more detail.


\subsection{Claim \ref{introclaim1} and the structure of singularities}
Claim \ref{introclaim1} is derived from a more localized statement, which
says that near points of large scalar curvature, the Ricci flow looks
very special : it is well-approximated by a special kind of
model Ricci flow, called a $\kappa$-solution.

\begin{proposition}[existence for kappa solutions]
  \label{prop:kl-overview-ricci-flow-approach-manifold-geometrization-existence-kappa-solutions}
  \dcref{Claim 7}
  \group{Kleiner--Lott Overview}
  \source{kleiner-lott}

\label{introclaim7}
Suppose that we have a given Ricci flow solution on a finite time interval.
If $x\in M$ and the scalar curvature $R(x,t)$ is large then
in the time-$t$ slice,
there is a ball  centered at $x$
of radius comparable to $R(x,t)^{-\frac12}$ in which the geometry
of the Ricci flow
is close to that of a ball in a $\kappa$-solution.
\end{proposition}
The quantity $R(x,t)^{-\frac12}$ is sometimes called the {\em curvature scale}
at $(x,t)$, because it scales like a distance.  We will  define $\kappa$-solutions
below, but mention here that they are Ricci flows with nonnegative sectional curvature,
and they are {\em ancient}, i.e. defined on a time interval of the form
$(-\infty,a)$.

The strength of Claim \ref{introclaim7} comes from the fact that there
is a good description of  $\kappa$-solutions.

\begin{proposition}[rigidity for kappa solutions]
  \label{prop:kl-overview-ricci-flow-approach-manifold-geometrization-rigidity-kappa-solutions}
  \dcref{Claim 8}
  \group{Kleiner--Lott Overview}
  \source{kleiner-lott}
 \cite{Perelman2} \label{introclaim8}
Any three-dimensional oriented $\kappa$-solution $(M_\infty, g_\infty(\cdot))$ falls
into one of the
following types : \\
(a) A finite isometric quotient of the round shrinking $3$-sphere. \\
(b) A Ricci flow on a  manifold  diffeomorphic to $S^3$ or $\R P^3$. \\
(c) A standard shrinking round neck on $\R \times S^2$ \\
(d) A Ricci flow on a manifold diffeomorphic to $\R^3$, each time slice of
which is  asymptotically necklike at infinity.\\
(e)  The  $\Z_2$-quotient $\R \times_{\Z_2} S^2$ of a  shrinking round neck.
\end{proposition}
Together, Claims \ref{introclaim7} and \ref{introclaim8} say that where the
scalar curvature is large, there is a region of diameter comparable to the
curvature scale  where one  sees either a closed manifold of known topology
(cases (a) and (b)),  a neck region
(case (c)),  a neck region capped off by a $3$-ball (case (d)), or a neck
region capped off by a twisted line bundle over $\R P^2$ (case (e)).
Applying this statement to every point of large scalar curvature at a time $t$
just prior to the blow-up time $T$,
one obtains a cover of $M$ by regions with special geometry and topology.
Any  overlaps occur in neck-like regions,
permitting one to splice them together to form the  connected  components
with known topology  whose existence is asserted in Claim \ref{introclaim1}.

Claim \ref{introclaim7} is proved using  a rescaling (or blow-up) argument.
This is a standard technique in geometric analysis and PDE's for treating
scale-invariant equations, such as the Ricci flow equation.
The claim is equivalent to the statement
that if $\{(x_i, t_i)\}_{i=1}^\infty$ is a sequence of spacetime points
for which  $\lim_{i \rightarrow \infty} R(x_i, t_i) = \infty$, then
by rescaling  the Ricci flow and passing to a subsequence, one obtains
a new sequence of Ricci flows which converges to a $\kappa$-solution.
More precisely,
view $(x_i, t_i)$ as a new
spacetime basepoint and spatially expand the solution around $(x_i, t_i)$
by $R(x_i,t_i)^{\frac12}$.  For dimensional reasons,
in order for rescaling to produce a new Ricci flow solution
one must also expand the time factor by $R(x_i,t_i)$. The new
Ricci flow solution, with time parameter $s$, is given by
\begin{equation}
\overline{g}_i(s) \: = \: R(x_i,t_i) \:
g \left( R(x_i,t_i)^{-1}  \: s \: + \: t_i \right).
\end{equation}
The new time interval for $s$ is
$\left[ - \:  R(x_i,t_i) \: t_i,0 \right]$.
One would then hope to take an appropriate limit $(M_\infty,
\overline{g}_\infty)$ of a subsequence of these rescaled solutions
$\{(M,  \overline{g}_i(\cdot))\}_{i=1}^\infty$.
(Technically speaking,
one uses smooth convergence of sequences of Ricci flows with basepoints;
this notion of convergence allows us to focus on what happens near
the spacetime points $(x_i, t_i)$.)
Any such  limit solution $(M_\infty,  \overline{g}_\infty)$ will be an ancient
solution,  since
$\lim_{i \rightarrow \infty} -R(x_i,t_i)\,t_i = -\infty$.
Furthermore, from a $3$-dimensional
result of Hamilton and Ivey (see Appendix \ref{phiappendix}),
any limit solution will have nonnegative sectional
curvature.


Although this sounds promising,
a major  problem was to show that a limit solution
actually exists.
To prove this, one would like to invoke Hamilton's compactness
theorem \cite{Hamiltonnn}.   In the present situation, the
compactness theorem says
that the sequence  of rescaled Ricci flows
$\{(M,  \overline{g}_i(\cdot))\}_{i=1}^\infty$
has a smoothly convergent subsequence provided two conditions are met: \\ \\
\noindent
A. For every $r>0$ and sufficiently large $i$,
the sectional curvature of $\overline{g_i}$ is  bounded uniformly independent of $i$
at each point $(x,s)$ in  spacetime such that $x$ lies in the $\overline{g_i}$-ball
$B_0(x_i,r)$ and $s\in [-r^2,0]$ and \\ \\
\noindent
B.  The  injectivity radii
$\inj(x_i, 0)$ in the time-$0$ slices of the $\overline{g_i}$'s have a
uniform positive lower bound.

For the moment, we ignore the issue of verifying condition A,
and simply assume that it holds for the sequence
$\{(M, \overline{g}_i(\cdot))\}_{i=1}^\infty$.
 In the presence
of the sectional curvature bounds in condition A,  a lower
bound on the injectivity
radius is known to be equivalent to a lower bound on the volume of
metric balls.  In terms of the original Ricci flow solution, this
becomes the condition that
\begin{equation} \label{bound}
r^{-3} \vol(B_t(x, r)) \ge \kappa \: > \: 0,
\end{equation}
where  $B_t(x, r)$ is
an arbitrary
metric $r$-ball in a time-$t$ slice, and the   curvature
bound   $|\Rm| \: \le \: \frac{1}{r^2}$ holds in $B_t(x,r)$.
The number $\kappa$ could depend on the
given Ricci flow solution, but the bound (\ref{bound})
should hold for all $t \in [0, T)$ and all $r < \rho$, where $\rho$ is a relevant
scale.


One of the outstanding achievements of \cite{Perelman} is to prove that
for an arbitrary Ricci flow defined on a
finite time interval, equation
(\ref{bound}) does hold with appropriate values of $\kappa$ and $\rho$.
In fact, the proof works in arbitrary dimension.
This result is called
a ``no local collapsing
theorem'' because it excludes the phenomenon of Cheeger-Gromov
collapse, in which a sequence of Riemannian manifolds has uniformly
bounded curvature, but fails to converge because the injectivity
radii tend to zero.

One can then
apply the no local collapsing theorem to the preceding rescaling argument,
provided that one has the needed sectional curvature bounds, in order to construct the
ancient solution $(M_\infty,\overline g_\infty)$.
In the blowup limit the condition that $r < \rho$ goes away, and so we
can say that $(M_\infty,\overline g_\infty)$ is $\kappa$-noncollapsed
(i.e. satisfies (\ref{bound})) at all scales. In addition, in the three-dimensional
case one can show that $(M_\infty,\overline g_\infty)$ has bounded sectional
curvature. To summarize, $(M_\infty,\overline g_\infty)$ is a
{\em $\kappa$-solution},
meaning that it is an ancient Ricci flow solution with nonnegative curvature operator on
each time slice and bounded sectional curvature
on compact time intervals,
which is $\kappa$-noncollapsed
at all scales.

With  the no local collapsing theorem in place, most of the proof of
Claim \ref{introclaim7} is concerned with showing that in the rescaling
argument, we effectively have the needed curvature
bounds of condition A.  The argument is a tour-de-force with many ingredients,
including earlier work
of Hamilton and the theory of
Riemannian manifolds with nonnegative sectional curvature.


\subsection{The proof of Claims \ref{introclaim2} and \ref{introclaim3}}
Claim \ref{introclaim1} allows one to take the limit of the evolving metric
$g(\cdot)$ as $t\ra T$, on the open set $\Omega$ where the metric is not
becoming singular.  It also provides geometrically defined regions -- the
connected components of the open set $U$ -- which one removes during surgery.
Each boundary component of the resulting manifold is
a nearly round $2$-sphere with a nearly cylindrical collar,
because the collar regions in Case B of Claim \ref{introclaim1} have a neck-like
geometry.  This enables one to glue in $3$-balls with a standard
metric, using a partition of unity  construction.


The Ricci flow starting with the postsurgery metric may also go singular after
a finite time. If so, one can appeal to Claim \ref{introclaim1} again to perform
surgery.
However, the elapsed time between successive surgeries will depend on
the scales at which surgeries are performed. Unless one performs
the surgeries very carefully, the surgery times may accumulate.

The way to rule out an accumulation of surgery times is to
arrange the surgery procedure so that a surgery at time $t$ removes a definite amount
of volume $v(t)$.  That is, a surgery at time $t$ should be performed at a
definite scale $h(t)$. In order to guarantee that this is possible, one needs to
establish a
quantitative version of Claim \ref{introclaim1} for a Ricci flow with surgery,
which applies not just at the first surgery time $T$ but also at a later
surgery
time $T^\prime$. The output of this quantitative version can depend on
the surgery time $T^\prime$
and the time-zero
metric, but it should be independent of whether or when surgeries occur
before time $T^\prime$.

The
general idea of the proof is similar to that of Claim \ref{introclaim1},
except that one has to carefully prescribe the surgery procedure in order
to control the effect of the earlier
surgeries. In particular, one of Perelman's
remarkable achievements is  a version of the no local collapsing theorem
for Ricci flows with surgery.


We refer the  reader to Section \ref{overview2} for a
more detailed overview of the proof of Claim \ref{introclaim3}, and for
further discussion of Claims \ref{introclaim5} and \ref{introclaim6}.



\section{Overview of
 {\em The Entropy Formula for the Ricci
Flow and its Geometric
Applications}
\cite{Perelman}}
\label{overview1}

The paper \cite{Perelman} deals with nonsingular Ricci flows;
the surgery process is considered in \cite{Perelman2}.
In particular, the final conclusion of \cite{Perelman} concerns
Ricci flows that are singularity-free and exist for all positive time.
It does not apply
to compact $3$-manifolds with finite fundamental group or
positive scalar curvature.

The purpose of the present overview is not to give a comprehensive
summary of the results of \cite{Perelman}. Rather we indicate its
organization and the interdependence of its sections, for the
convenience of the reader.
Some of the remarks in the overview
may only become clear once the reader has absorbed
a portion of the detailed notes.

Sections I.1-I.10, along with the first part of I.11, deal with Ricci flow on
$n$-dimensional manifolds.  The second part of I.11, and Sections I.12-I.13, deal
more specifically with Ricci flow on $3$-dimensional manifolds.
The main result is that geometrization holds if a compact $3$-manifold
admits a Riemannian metric which is the initial metric of a smooth
Ricci flow.  This was previously shown in \cite{Hamiltonn} under the
additional assumption that the sectional curvatures in the Ricci flow
are $O(t^{-1})$ as $t \rightarrow \infty$.

The paper \cite{Perelman} can be divided into four main parts.

Sections I.1-I.6 construct
certain entropy-type functionals ${\mathcal F}$ and ${\mathcal W}$
that are monotonic under Ricci flow. The functional ${\mathcal W}$ is
used to prove a no local collapsing theorem.

Sections I.7-I.10 introduce and apply another monotonic quantity,
the reduced volume $\tilde{V}$. It is also used to prove a
no local collapsing theorem.  The construction of $\tilde{V}$ uses
a modified notion of a geodesic, called an ${\mathcal L}$-geodesic.
For technical reasons the reduced volume $\tilde{V}$ seems to be
easier to work with than the ${\mathcal W}$-functional, and is used
in most of the sequel.
A reader who wants to focus on the Poincar\'e Conjecture or the
Geometrization Conjecture could
in principle start with I.7.

Section I.11 is concerned with $\kappa$-solutions, meaning nonflat
ancient solutions that are $\kappa$-noncollapsed at all scales
(for some $\kappa > 0$) and have bounded nonnegative curvature
operator on each time slice. In three dimensions, a blowup limit of a
finite-time singularity will be a $\kappa$-solution.

Sections I.12-I.13 are about three-dimensional Ricci flow solutions.
It is shown that high-scalar-curvature regions are modeled by
rescalings of $\kappa$-solutions. A decomposition of the
time-$t$ manifold into ``thick" and ``thin" pieces is described.
It is stated that as $t \rightarrow \infty$, the thick piece becomes
more and more hyperbolic, with incompressible cuspidal tori, and the thin
piece is a graph manifold. More details of these assertions
appear in \cite{Perelman2}, which also
deals with the necessary modifications if the solution has singularities.

We now describe each of these four parts in a bit more detail.

\subsection{I.1-I.6}

In these sections $M$ is assumed to be a closed $n$-dimensional
manifold.

A functional $F(g)$ of Riemannian metrics $g$ is said to be monotonic
under Ricci flow if $F(g(t))$ is nondecreasing in $t$ whenever $g(\cdot)$ is
a Ricci flow solution.
Monotonic quantities are an important tool for understanding Ricci flow. One wants
to have useful monotonic quantities, in particular with a characterization
of the Ricci flows for which $F(g(t))$ is constant in $t$.

Formally thinking of Ricci flow as a flow on the space of
metrics, one way to get a monotonic quantity  would be if the Ricci flow
were the gradient flow of a functional $F$. In Sections I.1-I.2, a functional ${\mathcal F}$
is introduced whose gradient flow is not quite Ricci flow, but only
differs from the Ricci flow by the action of diffeomorphisms.
(If one formally considers the Ricci flow as a flow on the space of
metrics modulo diffeomorphisms then it turns out to be the gradient flow of a functional
$\lambda_1$.) The functional ${\mathcal F}$ actually depends on a Riemannian metric
$g$ and a function $f$. If $g(\cdot)$ satisfies the Ricci flow equation
and $e^{-f(\cdot)}$ satisfies a conjugate or ``backward''
heat equation, in terms of $g(\cdot)$, then
${\mathcal F}(g(t),f(t))$ is nondecreasing in $t$. Furthermore, it is
constant in $t$ if and only if $g(\cdot)$ is a gradient steady soliton with
associated function $f(\cdot)$. Minimizing ${\mathcal F}(g,f)$ over all
functions $f$ with $\int_M e^{-f} \: dV \: = \: 1$ gives the monotonic
quantity $\lambda_1(g)$, which turns out to be the
lowest eigenvalue of $- \: 4 \: \triangle \: + \: R$.

In Section I.3 a modified ``entropy'' functional ${\mathcal W}(g,f,\tau)$ is
introduced.  It is nondecreasing in $t$ provided that $g(\cdot)$ is a Ricci flow,
$\tau \: = \: t_0 \: - \: t$ and $(4 \pi \tau)^{- \: \frac{n}{2}} \:
e^{-f}$ satisfies the conjugate heat equation. The functional
${\mathcal W}$ is constant on a Ricci flow if and only if the flow is a gradient
shrinking soliton that terminates at time $t_0$. Because of this,
${\mathcal W}$ is more suitable than ${\mathcal F}$ when one wants
information that is localized in spacetime.

In Section I.4 the entropy functional ${\mathcal W}$ is used to prove
a no local collapsing theorem. The statement is that if $g(\cdot)$ is a given
Ricci flow on a finite time interval $[0, T)$ then for any (scale) $\rho > 0$,
there is a number $\kappa > 0$ so that if $B_t(x,r)$ is a time-$t$ ball
with radius $r$ less than $\rho$, on which $|\Rm| \: \le \: \frac{1}{r^2}$,
then $\vol(B_t(x,r)) \: \ge \:
\kappa \: r^n$. The method of proof is to show that if
$r^{-n} \: \vol(B_t(x,r))$ is very small then the evaluation of ${\mathcal W}$
at time $t$ is very negative, which contradicts the monotonicity of
${\mathcal W}$.

The significance of a no local collapsing theorem is that it allows one to
use Hamilton's compactness theorem to construct blowup limits of finite time
singularities, and more generally to understand high curvature regions.

Section I.5 and I.6 are not needed in the sequel.  Section I.5 gives some
thermodynamic-like equations in which ${\mathcal W}$
appears as an entropy.  Section I.6 motivates the construction of the
reduced volume of Section I.7.

\subsection{I.7-I.10}

A new monotonic quantity, the reduced volume $\tilde{V}$,
is introduced in I.7.  It is defined in terms of so-called
${\mathcal L}$-geodesics. Let $(p, t_0)$ be a fixed spacetime point.
Define backward time by $\tau = t_0 -t$. Given a curve
$\gamma(\tau)$ in $M$ defined for $0 \le \tau \le \overline{\tau}$ (i.e.
going backward in real time) with $\gamma(0) \: = \: p$, its
${\mathcal L}$-length is
\begin{equation}
{\mathcal L}(\gamma) = \int_0^{\overline{\tau}} \sqrt{\tau} \:
\left( |\dot{\gamma}(\tau)|^2 \: + \: R(\gamma(\tau),t_0 - \tau) \right)
\: d\tau.
\end{equation}
One can compute the first and second variations of ${\mathcal L}$, in
analogy to what is done in Riemannian geometry.

Let $L(q, \overline{\tau})$
be the infimum of ${\mathcal L}(\gamma)$ over curves $\gamma$ with
$\gamma(0) = p$ and
$\gamma(\overline{\tau}) = q$. Put $l(q, \overline{\tau}) \: = \:
\frac{L(q, \overline{\tau})}{2\sqrt{\overline{\tau}}}$. The reduced volume
is defined by $\tilde{V}(\overline{\tau}) =
\int_M \overline{\tau}^{- \frac{n}{2}} \: e^{-l(q, \overline{\tau})} \:
\dvol(q)$. The remarkable fact is that if $g(\cdot)$ is a Ricci flow
solution then $\tilde{V}$ is nonincreasing in $\overline{\tau}$, i.e.
nondecreasing in real time $t$. Furthermore, it is constant if and only if
$g(\cdot)$ is a gradient shrinking soliton that terminates at time
$t_0$. The proof of monotonicity uses a subtle cancellation between
the $\overline{\tau}$-derivative of $l(\gamma(\overline{\tau}), \overline{\tau})$
along an ${\mathcal L}$-geodesic and the Jacobian of the
so-called ${\mathcal L}$-exponential map.

Using a differential inequality, it is shown that for each $\overline{\tau}$ there
is some point $q(\overline{\tau}) \in M$ so that
$l(q(\overline{\tau}), \overline{\tau}) \: \le \: \frac{n}{2}$. This is then
used to prove a no local collapsing theorem :  Given a Ricci flow
solution $g(\cdot)$
defined on a finite time interval $[0, t_0]$
and a scale $\rho > 0$ there
is a number $\kappa > 0$
with the following property. For $r < \rho$, suppose that
$|\Rm| \le \frac{1}{r^2}$ on the
``parabolic'' ball $\{(x,t) \: : \: \dist_{t_0}(x, p) \: \le \: r, \: \:
t_0 - r^2 \: \le \: t \: \le t_0\}$. Then
$\vol(B_{t_0}(p, r)) \: \ge \: \kappa \: r^n$. The number $\kappa$ can be chosen
to depend on $\rho$, $n$, $t_0$ and bounds on the geometry of the initial metric
$g(0)$.

The hypotheses of the no local collapsing theorem proved using $\tilde{V}$ are
more stringent than those of the no local collapsing theorem proved using $\mathcal{W}$,
but the consequences turn out to be the same.  The no local collapsing theorem is
used extensively in Sections I.11 and I.12 when extracting a convergent
subsequence from a sequence of Ricci flow solutions.

Theorem I.8.2 of Section I.8 says that under appropriate
assumptions, local $\kappa$-noncollapsing extends forwards in
time to larger distances.
This will be used in I.12 to analyze long-time behavior.
The statement is that for any $A < \infty$ there is a
$\kappa = \kappa(A) > 0$ with the following property. Given $r_0 > 0$ and
$x_0 \in M$, suppose that
$g(\cdot)$ is defined for $t \in [0, r_0^2]$, with
$|\Rm(x,t)| \: \le \: r_0^{-2}$ for all $(x,t)$ satisfying
$\dist_0(x, x_0) < r_0$, and the volume of the time-zero ball
$B_0(x_0, r_0)$ is at least $A^{-1} r_0^n$. Then the metric
$g(t)$ cannot be $\kappa$-collapsed on scales less than $r_0$
at a point $(x, r_0^2)$ with  $\dist_{r_0^2}(x, x_0) \le Ar_0$.

A localized version of the ${\mathcal W}$-functional appears in I.9.
Section I.10,
which is not needed for the sequel
but is of independent
interest, shows if a point in a time slice
lies in a ball with quantitatively bounded geometry then
at nearby later times, the curvature at the point is quantitatively
bounded.  That is, there is a damping effect with regard to exterior high-curvature
regions. The ``bounded geometry'' assumptions on the initial
ball are a lower bound on its scalar curvature and an assumption that
the isoperimetric constants of subballs are close to the Euclidean
value.

\subsection{I.11}

Section I.11 contains an analysis of $\kappa$-solutions. As mentioned before, in
three dimensions $\kappa$-solutions arise as blowup limits of finite-time singularities and,
more generally, as limits of rescalings of high-scalar-curvature regions.

In addition to the no local collapsing theorem, some of the tools used to analyze
$\kappa$-solutions are Hamilton's Harnack inequality for Ricci flows with
nonnegative curvature operator, and the comparison geometry of
nonnegatively curved manifolds.

The first result is that any time slice of a $\kappa$-solution
has vanishing asymptotic volume ratio $\lim_{r \rightarrow \infty} r^{-n} \vol(B_t(p, r))$.
This apparently technical result is used to show that if a
$\kappa$-solution $(M, g(\cdot))$ has scalar curvature normalized to equal
one at some spacetime point $(p,t)$ then there is an {\it a priori} upper bound on
the scalar curvature $R(q,t)$ at other points $q$ in terms of $\dist_t(p,q)$.
Using the curvature
bound, it is shown that a sequence $\{(M_i, p_i, g_i(\cdot)) \}_{i=1}^\infty$ of pointed
$n$-dimensional $\kappa$-solutions, normalized so that $R(p_i, 0) =1$ for each $i$, has
a convergent subsequence whose limit satisfies all of the requirements to be a
$\kappa$-solution, except possibly the condition of having bounded
sectional curvature on each time slice.

In three dimensions this statement is improved by showing that
the sectional curvature will be bounded on
each
compact time interval,
so the space of pointed $3$-dimensional $\kappa$-solutions
$(M, p, g(\cdot))$ with $R(p,0) = 1$ is in fact compact.
This is used to draw conclusions about the global geometry of
$3$-dimensional $\kappa$-solutions.

If $M$ is a compact $3$-dimensional $\kappa$-solution
then Hamilton's theorem about
compact $3$-manifolds with nonnegative Ricci curvature implies that $M$
is diffeomorphic to a spherical space form. If $M$ is noncompact then assuming
that $M$ is oriented, it follows easily that $M$ is
diffeomorphic to $\R^3$, or isometric to the round shrinking cylinder $\R \times S^2$
or its $\Z_2$-quotient $\R \times_{\Z_2} S^2$.

In the noncompact case it is shown that after rescaling, each time slice is neck-like at
infinity.  More precisely, considering a given time-$t$ slice, for each $\epsilon > 0$
there is a compact subset $M_\epsilon \subset M$ so that if $y \notin M_\epsilon$
then the pointed manifold $(M, y, R(y,t) g(t))$ is $\epsilon$-close
to the standard cylinder $\left[ - \: \frac{1}{\epsilon}, \frac{1}{\epsilon} \right] \times
S^2$ of scalar curvature one.

\subsection{I.12-I.13}

Sections I.12 and I.13 deal with three-dimensional Ricci flows.

Theorem I.12.1 uses the results of Section I.11 to model the
high-scalar-curvature regions of a Ricci flow.
Let us assume a
pinching condition of the form $\Rm \: \ge \: - \: \Phi(R)$ for an
appropriate function $\Phi$ with $\lim_{s \rightarrow \infty}
\frac{\Phi(s)}{s} = 0$.
(This will eventually follow from Hamilton-Ivey pinching,
 cf. Appendix \ref{phiappendix}.)
Theorem I.12.1 says that given
numbers $\epsilon, \kappa > 0$, one can find $r_0 > 0$ with the
following property. Suppose that $g(\cdot)$ is a Ricci flow solution
defined on some time interval $[0, T]$ that satisfies the
pinching condition and is $\kappa$-noncollapsed at scales less than
one.   Then for any point $(x_0, t_0)$ with $t_0 \ge 1$ and
$Q = R(x_0, t_0) \ge r_0^{-2}$, after scaling by the factor $Q$, the
solution in the region $\{(x,t) \: : \: \dist_{t_0}^2(x, x_0) \le (\epsilon Q)^{-1},
t_0 - (\epsilon Q)^{-1} \le t \le t_0\}$ is $\epsilon$-close to the
corresponding subset of a $\kappa$-solution.

Theorem I.12.1 says in particular that near a first singularity, the geometry is modeled
by a $\kappa$-solution, for some $\kappa$. This fact is used in
\cite{Perelman2}. Although Theorem I.12.1 is not
used directly in \cite{Perelman}, its method of proof is used in Theorem I.12.2.

The method of proof of Theorem I.12.1 is by contradiction. If it were not
true then there would be a sequence $r_0^{(i)} \rightarrow 0$ and
a sequence
$(M_i, g_i(\cdot))$ of
Ricci flow solutions that satisfy the assumptions, each with a spacetime point
$\left( x^{(i)}_0, t^{(i)}_0 \right)$ that does not satisfy the conclusion. To consider
first a special
case,
suppose that each point $\left( x^{(i)}_0, t^{(i)}_0 \right)$
is the first point at which a certain curvature threshold $R_i$ is
achieved, i.e.
$R(y, t) \: \le \: R \left( x^{(i)}_0, t^{(i)}_0 \right)$ for each
$y \in M_i$ and $t \in \left[ 0,  t^{(i)}_0 \right]$.
Then after rescaling the Ricci flow
$g_i(\cdot)$ by $Q_i \: = \: R \left( x^{(i)}_0, t^{(i)}_0 \right)$ and shifting the
time parameter, one has the curvature bounds on the time interval
$[- Q_i t^{(i)}_0,0]$ that form part of  the
hypotheses of Hamilton's compactness theorem.  Furthermore, the
no local collapsing theorem gives the lower injectivity radius bound
needed to apply Hamilton's theorem and take a convergent subsequence of the pointed
rescaled solutions.
The limit will be a $\kappa$-solution, giving the contradiction.

In the general case, one effectively proceeds by induction on the size of the
scalar curvature.  By modifying the choice of points $\left( x^{(i)}_0, t^{(i)}_0 \right)$,
one can assume that the conclusion of the theorem holds for all of the points $(y,t)$ in a
large spacetime neighborhood of $\left( x^{(i)}_0, t^{(i)}_0 \right)$
that have $R(y,t) > 2Q_i$. One then shows that one has the curvature bounds
needed to form the time-zero slice of the putative $\kappa$-solution. One shows
that this ``time-zero'' metric can be extended backward in time to form a
$\kappa$-solution, thereby giving the contradiction.

The rest of Section I.12 begins the analysis of the long-time behaviour of
a nonsingular $3$-dimensional Ricci flow. There are two main results,
Theorems I.12.2 and I.12.3. They extend curvature bounds forward and
backward in time, respectively.

Theorem I.12.2 roughly says that if
one has $|\Rm| \le r_0^{-2}$ on a spacetime region of spatial
size $r_0$ and temporal size $r_0^2$, and if one has a lower bound on the
volume of the initial time face of the region, then one gets scalar curvature
bounds on much larger spatial balls at the final time. More precisely,
for any $A < \infty$
there are numbers $K = K(A) < \infty$ and
$\rho = \rho(A) < \infty$ so that with the
hypotheses and notation of Theorem I.8.2, if in addition
$r_0^2 \: \Phi(r_0^{-2}) \: < \: \rho$
then
$R(x, r_0^2) \le K r_0^{-2}$ for points $x$ lying in the ball of
radius $Ar_0$
around $x_0$
at time $r_0^2$.

Theorem I.12.3 says that if one has a lower bound on volume and sectional
curvature on a ball at a certain time then one obtains an upper scalar
curvature bound on a smaller ball at an earlier time.  More precisely,
given $w > 0$ there exist $\tau = \tau(w) > 0$, $\rho = \rho(w) > 0$ and
$K = K(w) < \infty$ with the following property.  Suppose that
a ball
$B(x_0, r_0)$ at time $t_0$ has volume bounded below by
$w r_0^3$ and sectional curvature bounded below by $- \: r_0^{-2}$. Then
$R(x,t) < K r_0^{-2}$ for
$t \in [t_0 - \tau r_0^2, t_0]$ and $\dist_t(x,x_0) \: < \: \frac14 r_0$,
provided that $\phi(r_0^{-2}) < \rho$.

Applying a back-and-forth argument using Theorems I.12.2 and I.12.3,
along with the pinching condition, one concludes, roughly speaking,
that if a metric ball of small radius $r$ has infimal sectional curvature
exactly equal to $- \: r^{-2}$ then the ball has a small volume
compared to $r^3$. Such a ball can be said to be locally
volume-collapsed with respect to a lower sectional curvature bound.

Section I.13 defines the thick-thin decomposition of a large-time slice
of
a nonsingular
Ricci flow and shows the geometrization.
Rescaling the metric to $\widehat{g}(t) \: = \: t^{-1} \: g(t)$,
there is a universal function $\Phi$ so that for large $t$, the metric
$\widehat{g}(t)$ satisfies the $\Phi$-pinching condition. In terms of the
original unscaled metric, given $x \in M$ let $\widehat{r}(x,t) > 0$ be the unique
number such that $\inf \Rm \big|_{B_t(x,\widehat{r})} \: = \: - \:
\widehat{r}^{-2}$.

Given $w > 0$, define the $w$-thin part
$M_{thin}(w,t)$ of the time-$t$ slice to be the points $x \in M$ so that
$\vol(B_t(x, \widehat{r}(x,t))) \: < \: w \: \widehat{r}(x,t)^3$.
That is, a point
in $M_{thin}(w,t)$ lies in a ball that is locally volume-collapsed
with respect to a lower sectional curvature bound. Put
$M_{thick}(w,t) = M - M_{thin}(w,t)$. One shows that for large $t$, the subset
$M_{thick}(w,t)$ has bounded geometry in the sense that there
are numbers $\overline{\rho} = \overline{\rho}(w) > 0$ and
$K = K(w) < \infty$ so that $|\Rm| \: \le \: K t^{-1}$ on
$B(x, \overline{\rho} \sqrt{t})$ and
$\vol(B(x, \overline{\rho} \sqrt{t})) \: \ge \: \frac{1}{10} \: w \: (\overline{\rho} \sqrt{t}))^3$,
whenever $x \in M_{thick}(w,t)$.

Invoking arguments of Hamilton (that are written out in more detail in
\cite{Perelman2}) one can take a sequence $t \rightarrow \infty$ and
$w \rightarrow 0$ so that $M_{thick}(w,t)$ converges to a complete finite-volume
manifold with constant sectional curvature $- \: \frac14$, whose cuspidal
tori are incompressible in $M$. On the other hand, a result from Riemannian
geometry implies that for large $t$ and small $w$, $M_{thin}(w,t)$ is homeomorphic to
a graph manifold; again a more precise statement appears in \cite{Perelman2}.
The conclusion is that $M$ satisfies the geometrization conjecture. Again, one is
assuming in
I.13
that the Ricci flow is nonsingular
for all times.
